\documentclass[a4paper,10pt]{article}
\usepackage[top=0.5in,bottom=0.5in,left=0.55in,right=0.55in]{geometry}
\usepackage{enumitem}
\usepackage{titlesec}
\usepackage[T1]{fontenc}
\ifdefined\XeTeXrevision\else
  \usepackage[utf8]{inputenc}
\fi
\usepackage{lmodern}
\usepackage[hidelinks]{hyperref}
\usepackage{xcolor}
\ifdefined\XeTeXrevision\else
  \input{glyphtounicode}
  \pdfgentounicode=1
\fi

\makeatletter
% Allow \input files to resolve when LaTeX runs from repo root or cv/.
\def\input@path{{./}{cv/}}
\makeatother
\providecommand{\projectssectionfile}{sections/projects.es.tex}

\definecolor{linkblue}{HTML}{1a5276}
\hypersetup{colorlinks=true,urlcolor=linkblue}

\titleformat{\section}{\large\bfseries}{}{0em}{}[\titlerule]
\titlespacing{\section}{0pt}{8pt}{4pt}
\setlength{\parindent}{0pt}
\setlength{\parskip}{0pt}
\pagestyle{empty}

\begin{document}

\begin{center}
  \textbf{\LARGE Rody Vilchez}\\[4pt]
  IA Aplicada / ML Systems --- Retrieval \textperiodcentered{} Inteligencia Documental \textperiodcentered{} Evaluación y Robustez  \small
  \enspace\textbar\enspace
  \href{mailto:rody.vilchez00@gmail.com}{rody.vilchez00@gmail.com} 
  +51\,987\,082\,126 \enspace\textbar\enspace
  \href{https://rosewt.dev}{rosewt.dev} \enspace\textbar\enspace
  \href{https://github.com/R0SEWT}{GitHub} \enspace\textbar\enspace
  \href{https://www.linkedin.com/in/r0sewt/}{LinkedIn}
\end{center}

\vspace{2pt}

\small
Diseño sistemas de IA aplicada para condiciones no ideales: retrieval, inteligencia documental y pipelines de datos sobre corpus ruidosos y multilingües. Actualmente en CIP (CGIAR) construyendo pipelines de procesamiento documental y question answering sobre investigación agrícola. Publicación aceptada en Springer CCIS 2026 sobre traducción multimodal de lengua de señas. Mi foco está en la evaluación, la robustez y el comportamiento de modelos bajo restricciones reales.
\vspace{2pt}

\section*{Experiencia}

\textbf{AI / Data Intern} \hfill \textbf{International Potato Center (CIP, CGIAR)}\\
\textit{Lima, Perú} \hfill \textit{Oct 2025 -- Presente}
\begin{itemize}[leftmargin=1.2em, itemsep=1pt, topsep=2pt]
  \item Diseñé pipelines de procesamiento documental para alimentar un GraphRAG interno de CIP/CGIAR, sobre corpus multilingüe (español, inglés, francés, portugués, chino) con OCR ruidoso, layout irregular y clasificación parcial; ingesta, parsing, chunking, embedding y almacenamiento vectorial
  \item Implementé enriquecimiento de metadata con structured output basado en LLM (validación de esquema, batching, backoff frente a rate limits) para mejorar recuperación sobre documentos heterogéneos
  \item Diseñé y construí en conjunto un agente de soporte TI en Copilot Studio desplegado en Teams, cubriendo resolución de consultas nivel 0 (resolución sobre documentación técnica interna) y escalamiento a ticketing
  \item Diseñé el flujo de escalamiento: cuando el agente detecta que no puede resolver o el usuario lo pide explícitamente, genera un prellenado del ticket (structured output) a partir del contexto conversacional, con revisión human-in-the-loop vía Adaptive Cards antes del envío
  \item Rediseñé el agente de soporte TI, pasando de un diseño monolítico a un sistema multiagente hub-and-spoke, particionando las fuentes de conocimiento por dominio para mantenerse bajo el umbral de retrieval de la orquestación generativa de Copilot Studio y restaurar recuperación directa y sin filtrar por dominio
  \item Diseñé un framework de evaluación de tres capas, agnóstico de modelo, para agentes conversacionales: un runner determinista YAML-driven con análisis de matriz de confusión de ruteo, un scorer LLM-as-judge construido detrás de un contrato de modelo intercambiable (validado al migrar de proveedor de juez sin cambios en los consumidores), y tests de exploración con personas que detectan delegación silenciosa entre el agente padre y los hijos
  \item Extendí la plataforma interna de agentes de un único bot de soporte TI a una flota de agentes en Copilot Studio que abarca soporte TI, servicios administrativos corporativos y consultas sobre SOPs, construida sobre un submódulo de plataforma versionado con una ruta de upgrade de pin explícito y una plantilla de bootstrap que scaffoldea gobernanza, tests y tracking de issues para cada agente nuevo
  \item Construí un pipeline con arquitectura medallion (bronze/silver/gold) y un dashboard analítico sobre tickets históricos de soporte TI, segmentando usuarios mediante clustering no supervisado sobre embeddings y prototipando un reranker few-shot para clasificación de tickets basada en LLM con revisión human-in-the-loop
\end{itemize}

\vspace{2pt}

\textbf{QA Trainee} \hfill \textbf{Visma LATAM}\\
\textit{Lima, Perú} \hfill \textit{Dic 2024 -- Oct 2025}
\begin{itemize}[leftmargin=1.2em, itemsep=1pt, topsep=2pt]
  \item Construí un agente basado en LLM que genera tests automatizados end-to-end a partir de especificaciones, reduciendo el esfuerzo manual en la creación y mantenimiento de suites de regresión
  \item Desarrollé suites de regresión automatizadas con Cypress integradas en Jenkins, cubriendo flujos críticos que debían permanecer estables a través de integraciones sucesivas
  \item Construí generadores de pruebas DOM-aware para extraer selectores y estado desde aplicaciones en ejecución, mejorando mantenibilidad del testing frente a cambios de UI
\end{itemize}

\section*{Proyectos e Investigación}
\input{\projectssectionfile}

\section*{Educación}

\textbf{B.Sc.\ Ciencias de la Computación} \hfill \textbf{Universidad Peruana de Ciencias Aplicadas (UPC)}\\
\textit{Graduación estimada: 2026-2}

\section*{Reconocimientos y actividades}

DataFest - BCPxESAN (2do lugar, 2025) \textperiodcentered{}
SALA 2026 - Summit of AI in LatAm (full grant y participante, 2026) \textperiodcentered{}
Asociación KP (voluntariado, 95h, 2022--2023)

\section*{Habilidades}

\begin{tabular}{@{} l l @{} }
  \textbf{ML / AI Systems} & PyTorch, scikit-learn, Optuna, evaluación de modelos, pipelines multimodales \\
  \textbf{Retrieval / Document AI} & embeddings, Qdrant, LlamaIndex, chunking, parsing, procesamiento documental \\
  \textbf{Data / Backend} & Pandas, FastAPI, Flask, REST APIs, MongoDB, PostgreSQL, DuckDB, ETL \\
  \textbf{Infraestructura} & Docker, Git, Linux, Jenkins, GitHub Actions, Node.js, CI/CD \\
  \textbf{Plataformas de Agentes} & Copilot Studio, Power Automate, Dataverse, Adaptive Cards, Microsoft Teams, orquestación multiagente, evaluación LLM-as-judge \\
\end{tabular}

\vspace{4pt}

\section*{Certificaciones}

Developing Solutions for Microsoft Azure (AZ-204T00, WTC, 2026) \textperiodcentered{}
GitHub Foundations (GH-900T00, WTC, 2026) \textperiodcentered{}
AI Engineer for Data Scientists (DataCamp, 2025) \textperiodcentered{}
Machine Learning Specialization (Google Cloud, 2025) \textperiodcentered{}
Google Data Analytics (Google, 2024) \textperiodcentered{}
IA centrada en el ser humano (Tec de Monterrey, 2022)

\vspace{2pt}

\section*{Idiomas}

Español (nativo) \textperiodcentered{} Inglés (intermedio)

\end{document}
